%! AP Statistics — Unit 2, Lesson 3
\documentclass[10pt]{article}
\usepackage[margin=0.6in]{geometry}
\usepackage{xcolor}
\usepackage[most]{tcolorbox}
\usepackage{enumitem}
\usepackage{multicol}
\usepackage{amsmath,amssymb}
\usepackage{array}
\usepackage{tabularx}
\tcbuselibrary{breakable}

\definecolor{sky}{HTML}{EAF3FB}
\definecolor{navy}{HTML}{1F3A5F}
\definecolor{redbox}{HTML}{FCECEC}
\definecolor{greenbox}{HTML}{EDF8F0}
\definecolor{goldbox}{HTML}{FFF7E6}
\newtcolorbox{skillbox}[2][]{
    colback=#2, colframe=navy, boxrule=0.8pt, arc=2mm,
    left=2mm,right=2mm,top=1.5mm,bottom=1.5mm,
    fonttitle=\bfseries, title={#1}, halign title=center, breakable
}
\setlist[itemize]{leftmargin=1.2em,itemsep=0.35em,topsep=0.35em}
\setlist[enumerate]{leftmargin=1.4em,itemsep=0.35em,topsep=0.35em}
\newcolumntype{C}[1]{>{\centering\arraybackslash}p{#1}}
\newcolumntype{Y}{>{\centering\arraybackslash}X}

\newcommand{\CourseName}{AP Statistics}
\newcommand{\SchoolYear}{2026--2027}
\newcommand{\MeetingLength}{55 minutes}
\newcommand{\UnitNumberName}{Unit 2: Exploring Two-Variable Data \quad}
\newcommand{\LessonNumberName}{Lesson 2.3: Statistics for Two Categorical Variables}

\begin{document}
\begin{center}
    {\Large\bfseries \CourseName: \SchoolYear\ (\MeetingLength) \\
    {\small\bfseries \UnitNumberName \LessonNumberName}}
\end{center}

\begin{tcolorbox}[colback=sky,colframe=navy,boxrule=0.9pt,arc=2mm,
    left=2mm,right=2mm,top=1.5mm,bottom=1.5mm]
    \textbf{Primary Objective:} Students will quantify the association between two categorical
    variables by computing and interpreting the difference in proportions and relative risk.
    Students will use these statistics to make comparative claims in context, understanding that
    the size of the statistics (not just direction) conveys the strength of the association (Big Idea: UNC).
\end{tcolorbox}

\begin{skillbox}[Priority Ideas \& Skills]{goldbox}
\begin{minipage}[t]{0.48\linewidth}
    \textbf{AP Skill 2 --- Data Analysis}
    \begin{itemize}
        \item Compute conditional proportions (e.g., proportion of successes in each group).
        \item Calculate and interpret the difference in proportions in context.
        \item Calculate and interpret relative risk (ratio of proportions) in context.
        \item Recognize that a difference in proportions and relative risk can tell different stories about the same data.
    \end{itemize}
\end{minipage}
\hfill
\begin{minipage}[t]{0.48\linewidth}
    \textbf{Key Understandings}
    \begin{itemize}
        \item Difference in proportions is additive: ``Group A's rate is 0.15 higher than Group B's.''
        \item Relative risk is multiplicative: ``Group A's rate is 3 times Group B's rate.''
        \item Both statistics are needed: a small absolute difference can be a large relative risk (and vice versa).
        \item These statistics lay the groundwork for two-proportion inference in Unit 6.
    \end{itemize}
\end{minipage}
\end{skillbox}

\begin{skillbox}[{Vocabulary, Concepts, \& Theorems}]{greenbox}
\begin{minipage}[t]{0.48\linewidth}
    \begin{itemize}
        \item \textbf{Conditional proportion} --- the proportion of ``successes'' within a given group; $\hat{p} = \frac{\text{count of successes}}{\text{group total}}$
        \item \textbf{Difference in proportions} --- $\hat{p}_1 - \hat{p}_2$; quantifies the additive difference between two groups' rates
        \item \textbf{Relative risk (risk ratio)} --- $\dfrac{\hat{p}_1}{\hat{p}_2}$; how many times higher one group's rate is compared to the other
    \end{itemize}
\end{minipage}
\hfill
\begin{minipage}[t]{0.48\linewidth}
    \begin{itemize}
        \item \textbf{Baseline group} --- the reference group in the denominator of the relative risk
        \item \textbf{Strength of association} --- how much the conditional distributions differ; larger differences $\to$ stronger association
        \item \textbf{Direction of association} --- which group has the higher proportion; determines the sign of the difference
    \end{itemize}
\end{minipage}
\end{skillbox}

\begin{skillbox}[Activate Prior Knowledge \& Spiral Review]{sky}
\begin{minipage}[t]{0.48\linewidth}
    \textbf{Quick Review (3 min)}\\
    Return to the two-way table from Lesson 2.2. Students have already computed conditional relative frequencies. Ask: ``How much different are the two conditional proportions? Can we put a single number on that difference?'' This motivates both statistics today.
\end{minipage}
\hfill
\begin{minipage}[t]{0.48\linewidth}
    \textbf{Spiral Connections}
    \begin{itemize}
        \item Conditional relative frequencies (Lesson 2.2) are the $\hat{p}$ values used here.
        \item Unit 6: the difference in proportions becomes a test statistic for formal inference.
        \item Relative risk appears in epidemiology and medical research --- a real-world application of statistics.
    \end{itemize}
\end{minipage}
\end{skillbox}

\begin{skillbox}[Hook \& Lesson]{sky}
\begin{minipage}[t]{0.48\linewidth}
    \textbf{Hook (5 min)}\\
    Medical context: ``A study found that 2\% of people who took Drug A got a side effect, vs.\ 1\% who took a placebo.'' Ask: Is 1 percentage point a big deal? Now reframe: ``Those who took Drug A are \textit{twice} as likely to get the side effect.'' Discuss: same data, two different ways to describe it.
\end{minipage}
\hfill
\begin{minipage}[t]{0.48\linewidth}
    \textbf{Lesson Flow (15 min)}
    \begin{enumerate}
        \item Define and compute the difference in proportions from a two-way table.
        \item Define and compute relative risk from the same table.
        \item Interpret both in context side-by-side.
        \item Work a second example where relative risk is dramatic but absolute difference is tiny --- and vice versa.
        \item Discuss: which statistic is more informative depends on the question being asked.
    \end{enumerate}
\end{minipage}
\end{skillbox}

\begin{skillbox}[Explicit Instruction \& Active Monitoring]{sky}
\begin{minipage}[t]{0.48\linewidth}
    \textbf{Direct Instruction (10 min)}
    \begin{itemize}
        \item Difference in proportions template: ``The proportion of [response] among [Group 1] ($\hat{p}_1 = \ldots$) was [value] [higher/lower] than the proportion among [Group 2] ($\hat{p}_2 = \ldots$).''
        \item Relative risk template: ``[Group 1] was [value] times as likely to [response] as [Group 2].''
        \item Relative risk $= 1$ means no association; $> 1$ means Group 1 higher; $< 1$ means Group 1 lower.
    \end{itemize}
\end{minipage}
\hfill
\begin{minipage}[t]{0.48\linewidth}
    \textbf{Active Monitoring}
    \begin{itemize}
        \item Common error: interpreting relative risk as a percentage difference (e.g., ``RR = 2 means 2\% more'' --- wrong; it means twice as likely).
        \item Check: Is the student's interpretation in context? Is the direction stated?
        \item Probe: ``If the relative risk is 1.05, is there a strong association? What about the difference in proportions --- could it be large at the same time?''
    \end{itemize}
\end{minipage}
\end{skillbox}

\begin{skillbox}[Group Work \& Differentiation]{redbox}
\begin{minipage}[t]{0.48\linewidth}
    \textbf{Group Activity (8--10 min)}\\
    Groups receive a two-way table: vaccination status (vaccinated/unvaccinated) $\times$ illness (sick/not sick). Groups compute:
    \begin{enumerate}
        \item Conditional proportion sick for each group.
        \item Difference in proportions (with direction).
        \item Relative risk (vaccinated vs.\ unvaccinated).
        \item A 2-sentence conclusion about the association, including both statistics.
    \end{enumerate}
\end{minipage}
\hfill
\begin{minipage}[t]{0.48\linewidth}
    \textbf{Differentiation}
    \begin{itemize}
        \item \textit{Support}: Provide the interpretation templates filled in with blanks; students supply the numbers.
        \item \textit{Extension}: Students research and critique a real news article that reports relative risk without reporting the absolute difference --- is the headline misleading?
        \item \textit{AP Connection}: Discuss how Unit 6 tests will ask whether the observed difference in proportions is statistically significant.
    \end{itemize}
\end{minipage}
\end{skillbox}

\begin{skillbox}[Individual Work \& Assessment]{redbox}
\begin{minipage}[t]{0.48\linewidth}
    \textbf{Exit Ticket (5 min)}\\
    A survey finds 45 of 150 smokers (30\%) and 18 of 180 non-smokers (10\%) developed a certain condition. Students: (1) compute the difference in proportions and interpret, (2) compute the relative risk and interpret, (3) state which statistic better conveys the practical importance of the association and justify.
\end{minipage}
\hfill
\begin{minipage}[t]{0.48\linewidth}
    \textbf{AP-Style Check}\\
    \textit{``Based on the table, calculate and interpret the relative risk of [outcome] for Group A compared to Group B. Does the relative risk or the difference in proportions better communicate the strength of this association? Explain.''}\\[4pt]
    Assess: correct calculation, correct direction, full context in interpretation.
\end{minipage}
\end{skillbox}

\begin{skillbox}[Reinforcement \& Extension]{goldbox}
\begin{minipage}[t]{0.48\linewidth}
    \textbf{Homework / Practice}
    \begin{itemize}
        \item Given two different two-way tables, compute both statistics for each and write a comparison paragraph.
        \item Reflect: When would a researcher prefer to report relative risk vs.\ difference in proportions?
    \end{itemize}
\end{minipage}
\hfill
\begin{minipage}[t]{0.48\linewidth}
    \textbf{Extension / Preview}
    \begin{itemize}
        \item \textit{Extension}: Research ``absolute risk reduction'' and ``number needed to treat'' --- two more ways to quantify the same association in medical studies.
        \item \textit{Preview}: Lesson 2.4 shifts to two \textit{quantitative} variables. Instead of a table, we'll plot data on a coordinate plane --- a scatterplot. Think about what features you would look for in such a plot.
    \end{itemize}
\end{minipage}
\end{skillbox}

\end{document}
